\subsection{Mechanism and failure traces}
\textbf{Representation and evidence (244).}
In the old C+ reply, the model cited a transaction identifier as if it were a receipt identifier. In the V5 revision the cited receipt still failed to establish a completed step, but the usable intent and independent edits remained. The executor received a partial state. A successfully parsed decision to withdraw an action did not automatically establish completion of the affected step.

\textbf{Guard control and argument coverage (750 and 737).}
On the separate pilot checkpoint (750), the user kept an attachment and revoked sending a draft. A controlled invocation of the proposed send returned HTTP 409 with no simulator submission and identical state fingerprints. A subsequent lookup returned HTTP 200. The formal return episode (737) illustrates a different outcome. G5 made a rule for the original pending return, but six later requests used different order identities. The guard allowed them. All six returned HTTP 200 and added return records under the monitor. The narrow identity constraint missed those application writes.

\textbf{Notification visibility (267 and 272).}
The resumed branches reused a payment request that had already been sent. The standalone branches made the requested payment request from a fresh world. All saved terminal comparisons yielded one added notification per branch under explicit inclusion. The upstream evaluator's auxiliary ignore list removed that model from a test which required it. The compatibility score restores the notification check for every branch without changing any action logs. This resolves an evaluator inconsistency separately from the performance contrast.
